%% toc: P15 173 173.9,173.10,173.11,173.12
\section{Proofs for computed and simulated values in approximate conclusions}
\label{sec:annex-p15-w02b-p15-074}

Notation is that of \cref{def:p15-w02b-p15-074-ledger}. In decimal
precision three, \(R\) is the nearest rounding of
\cref{def:p12-numerical-rounding}; for a real \(z\) with
\(10^e\leq|z|<10^{e+1}\), \(\operatorname{ulp}(z)=10^{e-2}\), and
\cref{thm:p12-numerical-nearest} gives \(|R(z)-z|\leq\tfrac12\cdot10^{e-2}\).

\begin{annexitem}{annex:p15-w02b-p15-074-ledger}{thm:p15-w02b-p15-074-ledger}{The telescoped ledger, its certificates and its sharpness}
\begin{proof}
Write \(d_1=Y-F(x_*)\), \(d_2=F(x_*)-A(x_*)\),
\(d_3=A(x_*)-A(\widehat x)\) and \(d_4=A(\widehat x)-\widehat A\), so
that \(Y-\widehat A=d_1+d_2+d_3+d_4\). The model premise gives
\(-E_{\rm model}\leq d_1\leq E_{\rm model}\). Since \(x_*\in B\), the
approximation premise at \(x=x_*\) gives
\(\ell_{\rm appr}\leq d_2\leq r_{\rm appr}\), and the input premise at
\(x=x_*\) gives \(-E_{\rm input}\leq d_3\leq E_{\rm input}\). The
computation premise gives \(-E_{\rm comp}\leq d_4\leq E_{\rm comp}\).
Adding the four lower bounds and the four upper bounds yields
\(\ell_{\rm appr}-E\leq Y-\widehat A\leq r_{\rm appr}+E\), which is the
displayed interval.

Item 1. If \(\widehat A+\ell_{\rm appr}-E>c\), then
\(Y\geq\widehat A+\ell_{\rm appr}-E>c\). If
\(\widehat A+r_{\rm appr}+E<c\), then
\(Y\leq\widehat A+r_{\rm appr}+E<c\).

Item 2. Applying the interval to each target,
\(Y_1\leq\widehat A_1+r_1+E_1<\widehat A_2+\ell_2-E_2\leq Y_2\).

Item 3. Let \(J\) denote the displayed interval. Every value of \(Y\)
compatible with the premises lies in \(J\) by the first paragraph.
Conversely let \(y\in J\). The map
\((d_1,d_2,d_3,d_4)\mapsto d_1+d_2+d_3+d_4\) is continuous on the box
\([-E_{\rm model},E_{\rm model}]\times[\ell_{\rm appr},r_{\rm appr}]
\times[-E_{\rm input},E_{\rm input}]\times[-E_{\rm comp},E_{\rm comp}]\),
which is connected; its values at the lowest and highest corners are
\(\ell_{\rm appr}-E\) and \(r_{\rm appr}+E\). By the intermediate value
theorem along the segment joining these corners, some point
\((d_1,d_2,d_3,d_4)\) of the box has sum \(y-\widehat A\). Put
\(A(\widehat x)=\widehat A+d_4\), \(A(x_*)=A(\widehat x)+d_3\),
\(F(x_*)=A(x_*)+d_2\) and \(Y=F(x_*)+d_1=y\). These numbers satisfy the
four premises at \(x=x_*\). They are also realized by functions: when
\(d_3\neq0\), take \(S=B=\{x_*,\widehat x\}\) with two distinct points
and set \(F(\widehat x)=A(\widehat x)+d_2\); then the approximation premise
holds at both points, the input premise holds at \(\widehat x\) trivially
and at \(x_*\) because \(|d_3|\leq E_{\rm input}\). When \(d_3=0\), take
\(S=B=\{x_*\}\) and \(\widehat x=x_*\). So every \(y\in J\) is a
compatible value. If \(c\in J\), the compatible assignment with \(Y=c\)
refutes both \(Y>c\) and \(Y<c\), so neither follows from the premises.

Item 4. In the symmetric case write the four nonnegative magnitudes as
\(E_1,\ldots,E_4\). Then
\[
 \Bigl(\sum_iE_i\Bigr)^2=\sum_iE_i^2+2\sum_{i<j}E_iE_j,
\]
and the cross sum is positive when at least two magnitudes are positive.
Hence \(\bigl(\sum_iE_i^2\bigr)^{1/2}<\sum_iE_i\). The assignment of
item 3 with every \(d_k\) at its upper endpoint satisfies the premises and
gives \(|Y-\widehat A|=\sum_iE_i\), which exceeds the root-sum-square.
\end{proof}
\end{annexitem}

\begin{annexitem}{annex:p15-w02b-p15-074-regime}{thm:p15-w02b-p15-074-regime}{Enclosing a regime parameter}
\begin{proof}
Let \(x\in B\). By the input premise for \(\varphi\), the computation
premise and the hypothesis,
\[
 \varphi(x)\leq\varphi(\widehat x)+H_{\rm input}
 \leq\widehat\varphi+H_{\rm arith}+H_{\rm input}\leq\rho,
\]
so \(x\in D\). Hence \(B\subseteq D\), and a remainder bound proved for
every point of \(D\) holds at every point of \(B\), which is what the
approximation entry of \cref{def:p15-w02b-p15-074-ledger} requires.

For the second claim take \(S=\mathbb R\), \(\varphi(x)=x\),
\(\rho=1/2\), \(\widehat x=1/2\), \(\widehat\varphi=1/2\) and, for some
\(\eta>0\), \(B=[1/2-\eta,1/2+\eta]\). The premises hold with
\(H_{\rm arith}=0\) and \(H_{\rm input}=\eta\), and
\(\widehat\varphi=\varphi(\widehat x)\leq\rho\). Yet
\(1/2+\eta\in B\) and \(\varphi(1/2+\eta)>\rho\), so
\(B\not\subseteq D\). This is the domain restriction of
\cref{thm:p15-approximation-rational}.
\end{proof}
\end{annexitem}

\begin{annexitem}{annex:p15-w02b-p15-074-four-thresholds}{thm:p15-w02b-p15-074-four-thresholds}{Four threshold calculations}
\begin{proof}
Item 1. The operands \(1\) and \(0.123=1.23\times10^{-1}\) are values of
the format, so the single rounding is that of the sum
\(z=1.123\). Since \(1\leq z<10\), \(\operatorname{ulp}(z)=0.01\); the
neighbouring format values are \(1.12\) and \(1.13\), at distances
\(0.003\) and \(0.007\), so \(\widehat A=R(z)=1.12\). The half-ulp bound
gives \(|A(\widehat x)-\widehat A|\leq0.005\), so \(E_{\rm comp}=0.005\)
is admissible. Every \(x\in B\) satisfies \(0\leq x\leq1/2\). By
\cref{thm:p15-approximation-rational},
\(0\leq F(x)-A(x)\leq2x^2\leq2(0.125)^2=1/32\) on \(B\), giving
\(\ell_{\rm appr}=0\) and \(r_{\rm appr}=1/32\). For \(x\in B\),
\(|A(x)-A(\widehat x)|=|x-0.123|\leq0.002\). With \(E_{\rm model}=0\),
\(E=0.007\), and \cref{thm:p15-w02b-p15-074-ledger} gives the interval
\([1.12-0.007,\,1.12+0.03125+0.007]=[1.113,\,1.15825]\).

Item 2. \(1.113>1.10\), so item 1 of
\cref{thm:p15-w02b-p15-074-ledger} certifies \(Y>1.10\).

Item 3. \(1.113\leq1.12\leq1.15825\), so by item 3 of
\cref{thm:p15-w02b-p15-074-ledger} the ledger certifies neither strict
comparison. The exact value \(A(\widehat x)=1.123\) needs no rounding
analysis, so \(E_{\rm comp}=0\) and the interval is
\([1.123-0.002,\,1.123+0.03125+0.002]=[1.121,\,1.15625]\). Its lower
endpoint exceeds \(1.12\).

Item 4. \(1.121\leq1.13\leq1.15625\). For \(0\leq x<y<1\),
\(1-x>1-y>0\), so \(F(x)<F(y)\): \(F\) is increasing on \(B\). Hence
\(F(x_*)\in[F(0.121),F(0.125)]=[1000/879,\,8/7]\). Since
\(1.137\cdot879=999.423<1000\), \(1000/879>1.137>1.13\), and therefore
\(Y=F(x_*)>1.13\). No rounding occurs: the endpoints are exact rationals.

Item 5. \(1.14\cdot879=1002.06>1000\) gives \(F(0.121)=1000/879<1.14\),
and \(8/7>1.14\) because \(8>7.98\). Both \(0.121\) and \(0.125\) satisfy
\(|x-0.123|\leq0.002\). With \(Y=F(x_*)\) computed exactly, the choice
\(x_*=0.121\) gives \(Y<1.14\) and the choice \(x_*=0.125\) gives
\(Y>1.14\). Both are compatible with every premise, and neither the
arithmetic nor the approximation entry constrains \(x_*\); they are
already zero in this calculation. So no improvement of those entries
decides the comparison.

Item 6. With a model entry, the generic lower endpoint becomes
\(1.113-E_{\rm model}\), which exceeds \(1.10\) if and only if
\(E_{\rm model}<0.013\). In item 4 the lower bound becomes
\(1000/879-E_{\rm model}\), which exceeds \(1.13\) if and only if
\[
 E_{\rm model}<\frac{1000}{879}-\frac{113}{100}
 =\frac{100000-99327}{87900}=\frac{673}{87900}.
\]
\end{proof}
\end{annexitem}

\begin{annexitem}{annex:p15-w02b-p15-074-euler}{thm:p15-w02b-p15-074-euler}{Euler thresholds}
\begin{proof}
Throughout \(q_0=\lambda=T=1\) in \cref{def:p12-simulation-decay}.

Item 1. Exactly,
\[
 v_{20}=\frac{19^{20}}{20^{20}}
 =\frac{37589973457545958193355601}{104857600000000000000000000}.
\]
Multiplying the denominator by \(0.35848\) gives
\(37589352448000000000000000\), which is smaller than the numerator, and
multiplying it by \(0.35849\) gives \(37590401024000000000000000\), which is
larger. So \(0.35848<v_{20}<0.35849\). The step \(h=1/20\) satisfies
\(0<\lambda h<1\), so \cref{thm:p12-simulation-decay-euler} gives
\[
 0\leq q(1)-v_{20}\leq q(1)\frac{1/20}{2(19/20)}=\frac{q(1)}{38}.
\]
The left inequality is \(v_{20}\leq q(1)\); the right one rearranges to
\(\tfrac{37}{38}q(1)\leq v_{20}\), that is \(q(1)\leq\tfrac{38}{37}v_{20}\).
Finally \(\tfrac{38}{37}v_{20}<\tfrac{38}{37}\cdot0.35849
=13.62262/37<0.36818\), because \(0.36818\cdot37=13.62266\).

Item 2. The exact products and their roundings are, in order,
\(0.95\mapsto0.95\), \(0.9025\mapsto0.902\), \(0.8569\mapsto0.857\),
\(0.81415\mapsto0.814\), \(0.7733\mapsto0.773\), \(0.73435\mapsto0.734\),
\(0.6973\mapsto0.697\), \(0.66215\mapsto0.662\), \(0.6289\mapsto0.629\),
\(0.59755\mapsto0.598\), \(0.5681\mapsto0.568\), \(0.5396\mapsto0.540\),
\(0.513\mapsto0.513\), \(0.48735\mapsto0.487\), \(0.46265\mapsto0.463\),
\(0.43985\mapsto0.440\), \(0.418\mapsto0.418\), \(0.3971\mapsto0.397\),
\(0.37715\mapsto0.377\), \(0.35815\mapsto0.358\).
The second step is an exact midpoint, resolved to the even last digit as
\cref{def:p12-numerical-rounding} prescribes; every other step has a unique
nearest value. So \(\widehat v_{20}=0.358\). Every exact product lies in
\([0.1,1)\), where the ulp is \(0.001\), so the rounding error
\(\delta_j=\widehat v_{j+1}-0.95\,\widehat v_j\) satisfies
\(|\delta_j|\leq0.0005\). The coefficient \(0.95=1-a\) with \(a=1/20\) is a
format value and is used exactly, and \(\widehat v_0=v_0=1\). Thus
\cref{thm:p12-simulation-arithmetic-propagation} applies with \(E_0=0\) and
\(\rho=0.0005\):
\[
\begin{aligned}
 |\widehat v_{20}-v_{20}|&\leq\frac{0.0005\,(1-0.95^{20})}{0.05}
 =0.01\,(1-v_{20})\\
 &<0.01\,(1-0.35848)=0.0064152<0.0065.
\end{aligned}
\]

Item 3. By items 1 and 2,
\(q(1)\geq v_{20}\geq\widehat v_{20}-0.01(1-v_{20})>0.358-0.0064152
=0.3515848>0.35\). For \(0.355\), the lower endpoint this budget supplies
is \(\widehat v_{20}-0.01(1-v_{20})<0.358-0.01(1-0.35849)=0.3515849<0.355\),
while its upper endpoint \(\tfrac{38}{37}\bigl(\widehat v_{20}
+0.01(1-v_{20})\bigr)\) exceeds \(0.358>0.355\). The interval contains
\(0.355\), so by \cref{thm:p15-approximation-certificates} it certifies
neither strict comparison. The exact-real iterate gives
\(q(1)\geq v_{20}>0.35848>0.355\).

Item 4. By item 1, \(q(1)\in[v_{20},\tfrac{38}{37}v_{20}]\), with
\(v_{20}<0.35849<0.365\) and
\(\tfrac{38}{37}v_{20}>\tfrac{38}{37}\cdot0.35848>0.3681>0.365\).
The interval contains \(0.365\), so by
\cref{thm:p15-approximation-certificates} it certifies neither strict
comparison. For \(N=100\), \(h=1/100\) satisfies \(0<\lambda h<1\), and
\cref{thm:p12-simulation-decay-euler} gives \(q(1)\geq v_{100}=0.99^{100}\).
For \(0\leq x<1\), the series \(\log(1-x)=-\sum_{k\geq1}x^k/k\) gives
\[
 \log(1-x)\geq-x-\frac{x^2}{2}\sum_{k\geq0}x^k=-x-\frac{x^2}{2(1-x)}.
\]
With \(x=1/100\), \(100\log0.99\geq-1-1/198\). Hence
\(v_{100}\geq e^{-1}e^{-1/198}\geq e^{-1}(1-1/198)\), using
\(e^{-y}\geq1-y\). In the series \(e^{-1}=\sum_{k\geq0}(-1)^k/k!\), the
terms after \(k=7\) group into the positive pairs
\(1/(2m)!-1/(2m+1)!\) for \(m\geq4\), so \(e^{-1}\) exceeds the partial
sum through \(k=7\):
\[
 1-1+\tfrac12-\tfrac16+\tfrac1{24}-\tfrac1{120}+\tfrac1{720}-\tfrac1{5040}
 =\frac{1854}{5040}.
\]
Therefore \(v_{100}\geq\tfrac{1854}{5040}\cdot\tfrac{197}{198}
=\tfrac{365238}{997920}\), and \(0.365\cdot997920=364240.8<365238\), so
\(v_{100}>0.365\) and \(q(1)>0.365\).
\end{proof}
\end{annexitem}

\begin{annexitem}{annex:p15-w02b-p15-074-sampled}{thm:p15-w02b-p15-074-sampled}{Indicator band, sampling bound and decision error}
\begin{proof}
Item 1. Since \(G\), \(\widehat G\) and \(\xi\) are measurable, the
indicators \(I=\mathbf 1\{G(\xi)\geq a\}\) and
\(\widehat I=\mathbf 1\{\widehat G(\xi)\geq a\}\) are random variables.
If \(G(s)-a>E\), then \(\widehat G(s)\geq G(s)-E>a\), and both indicators
equal one at \(s\). If \(G(s)-a<-E\), then \(\widehat G(s)\leq G(s)+E<a\),
and both equal zero. So \(\{\widehat I\neq I\}\subseteq
\{|G(\xi)-a|\leq E\}\), and
\[
 |\widetilde p-p|=\bigl|\mathbb E[\widehat I-I]\bigr|
 \leq\mathbb E|\widehat I-I|=\mathbb P(\widehat I\neq I)\leq\omega(E).
\]

Item 2. Put \(Z_i=\mathbf 1\{\widehat G(\xi_i)\geq a\}\). Each \(Z_i\) is
a measurable function of \(\xi_i\), so the \(Z_i\) are independent by
\cref{thm:p15-independent-blocks}, and each has the law of \(\widehat I\):
\(\mathbb E Z_i=\widetilde p\). Since \(Z_i^2=Z_i\),
\(\operatorname{Var}(Z_i)=\widetilde p-\widetilde p^{\,2}
=\tfrac14-(\widetilde p-\tfrac12)^2\leq\tfrac14\). Independence gives zero
covariances, so by \cref{thm:p15-independent-moments}
\(\operatorname{Var}(\widehat p_n)=n^{-2}\sum_i\operatorname{Var}(Z_i)
\leq1/(4n)\), while \(\mathbb E\widehat p_n=\widetilde p\). The variance
form of Chebyshev's bound in \cref{thm:p15-expectation-bounds} gives
\(\mathbb P(|\widehat p_n-\widetilde p|\geq t)\leq1/(4nt^2)\).

Item 3. Suppose \(p\leq p_0\). By item 1,
\(\widetilde p\leq p+\omega(E)\leq p_0+\omega(E)\). On the event
\(\widehat p_n>p_0+t+\omega(E)\), therefore,
\(\widehat p_n-\widetilde p>t\). That event lies in
\(\{|\widehat p_n-\widetilde p|\geq t\}\), whose probability is at most
\(1/(4nt^2)\) by item 2. If \(p\geq p_0\), then
\(\widetilde p\geq p_0-\omega(E)\), and on the event
\(\widehat p_n<p_0-t-\omega(E)\) we have \(\widetilde p-\widehat p_n>t\);
the same bound applies.
\end{proof}
\end{annexitem}

\begin{annexitem}{annex:p15-w02b-p15-074-tie}{thm:p15-w02b-p15-074-tie}{The decimal tie}
\begin{proof}
Take \(S=\{1,2,3\}\) with all subsets measurable. In the reals
\(G(s)=s\), so \(G(\xi)\geq1\) always and \(p=1\).

For \(s=1\): \(1/3\) lies in \([0.1,1)\), where the ulp is \(0.001\); its
neighbours \(0.333\) and \(0.334\) are at distances \(1/3000\) and
\(2/3000\), so \(R(1/3)=0.333\). The product \(0.333\cdot3=0.999\) is a
format value, so \(\widehat G(1)=0.999\). For \(s=2\): the neighbours of
\(2/3\) are \(0.666\) and \(0.667\), at distances \(2/3000\) and
\(1/3000\), so \(R(2/3)=0.667\); then \(0.667\cdot3=2.001\) lies in
\([1,10)\), where the ulp is \(0.01\), and its neighbours \(2.00\) and
\(2.01\) are at distances \(0.001\) and \(0.009\), so \(\widehat G(2)=2\).
For \(s=3\) both operations are exact and \(\widehat G(3)=3\).
The errors are \(0.001\), \(0\) and \(0\), so \(E=0.001\) is a uniform
bound. Only \(s=1\) fails the computed test, so
\(\widetilde p=\mathbb P(\xi\in\{2,3\})=2/3\). Since \(|G(s)-1|\leq0.001\)
holds only at \(s=1\), \(\omega(E)=1/3=|\widetilde p-p|\).

The rule of item 3 declares \(p>0.9\) only if
\(\widehat p_n>0.9+t+1/3>1\), which is impossible for a frequency.
Since \(p=1\geq0.9\), item 3 bounds the probability that it declares
\(p<0.9\) by \(1/(4nt^2)\).

Without \(\omega(E)\), the rule declares \(p<0.9\) when
\(\widehat p_n<0.9-t\). For \(0<t<7/30\), the complementary event
\(\widehat p_n\geq0.9-t\) implies
\(\widehat p_n-\widetilde p\geq0.9-t-2/3=7/30-t>0\), whose probability is
at most \(1/\bigl(4n(7/30-t)^2\bigr)\) by item 2. The false declaration
therefore has probability at least \(1-1/\bigl(4n(7/30-t)^2\bigr)\).
\end{proof}
\end{annexitem}
